\section*{अध्याय \ref{ch:g6:water-underground} --- ज़मीन के नीचे का पानी}

\begin{solution}{exo:g6:water-underground:1}
\omterm{def:g6:water-underground:permeability}{पारगम्य}: \omterm{def:g3:stones-pebbles-sand:sand}{रेत}, \omterm{def:g3:stones-pebbles-sand:pebble}{बजरी}, खड़िया। \omterm{def:g6:water-underground:permeability}{अपारगम्य}: चिकनी मिट्टी, बिना दरारों वाला ग्रेनाइट,
शेल।
\end{solution}

\begin{solution}{exo:g6:water-underground:2}
(क) \omterm{def:g6:water-underground:water-table}{भौम जलस्तर}; (ख) \omterm{def:g6:water-underground:aquifer}{जलभृत}; (ग) \omterm{def:g6:water-underground:spring}{सोता}; (घ) \omterm{def:g6:water-underground:spring}{कुआँ}।
\end{solution}

\begin{solution}{exo:g6:water-underground:3}
(क) वह वाष्पित होकर वापस हवा में चला जाता है; वह सतह के ऊपर से बहकर नालों में
चला जाता है; वह ज़मीन में रिस जाता है (\omterm{def:g6:water-underground:infiltration}{अंतःस्यंदन})। (ख) कार पार्किंग पर: उसकी
सख़्त सतह लगभग कोई पानी रिसने नहीं देती, इसलिए पानी बह जाता है। घास के मैदान में
मिट्टी और घास उसके बड़े हिस्से को रिसने देती हैं।
\end{solution}

\begin{solution}{exo:g6:water-underground:4}
(क) ग़लत: चिकनी मिट्टी बहुत \omterm{def:g6:water-underground:porosity}{सरंध्र} है पर \omterm{def:g6:water-underground:permeability}{अपारगम्य}, क्योंकि उसके \omterm{def:g6:water-underground:porosity}{रंध्र} नन्हे हैं।
(ख) सही। (ग) ग़लत: वह ज़मीन के आकार के साथ चलता है, उससे कम उतार-चढ़ाव के साथ,
और मौसमों के साथ ऊपर-नीचे होता है। (घ) सही: सारे \omterm{def:g2:water-on-earth:salt-fresh}{मीठे पानी} का लगभग 30~\%
\omterm{def:g6:water-underground:groundwater}{भूजल} है, जबकि \omterm{def:g1:landscapes-around-us:water}{झीलों} में लगभग एक प्रतिशत का चौथाई हिस्सा है।
\end{solution}

\begin{solution}{exo:g6:water-underground:5}
$10 \times 25 \div 100 = 2.5$: बाल्टी के उफनने से पहले तुम 2.5~लीटर डाल सकते हो।
$10 \times 22 \div 100 = 2.2$: छेद से 2.2~लीटर बाहर निकलेगा।
\end{solution}

\begin{solution}{exo:g6:water-underground:6}
(क) $2 \times 1 \times 1 = 2$~घन मीटर, यानी 2000~लीटर।
(ख) $2000 \times 20 \div 100 = 400$~लीटर। (ग) $2000 \times 18 \div 100 =
360$~लीटर बाहर निकलते हैं; $400 - 360 = 40$~लीटर अटके रहते हैं।
\end{solution}

\begin{solution}{exo:g6:water-underground:7}
(क) \omterm{def:g6:water-underground:spring}{सोता} लगभग 203~मीटर पर है, \omterm{def:g1:landscapes-around-us:water}{नदी} लगभग 132~मीटर पर। (ख) \omterm{def:g6:water-underground:spring}{सोता} \omterm{def:g1:landscapes-around-us:water}{नदी} से लगभग
70~मीटर ऊपर है: उसका पानी ढलान से नीचे बहकर \omterm{def:g1:landscapes-around-us:water}{नदी} में जाता है। (ग) खड़िया की
\omterm{def:g1:landscapes-around-us:landform}{पहाड़ी} ढलान के अंदर \omterm{def:g6:water-underground:water-table}{भौम जलस्तर} ज़मीन के नीचे है; वह सतह से केवल वहीं मिलता है
जहाँ खड़िया का आधार, जो \omterm{def:g6:water-underground:permeability}{अपारगम्य} शेल पर टिका है, ढलान पर बाहर
निकलता है।
\end{solution}

\begin{solution}{exo:g6:water-underground:8}
(क) बरसात के मौसम के बाद: $401 - 336 = 65$~मीटर; सूखे मौसम के बाद:
$401 - 314 = 87$~मीटर। (ख) $336 - 314 = 22$~मीटर।
\end{solution}

\begin{solution}{exo:g6:water-underground:9}
ज़्यादा गहरा। \omterm{def:g6:water-underground:water-table}{भौम जलस्तर} ज़मीन के आकार के साथ चलता है, पर उससे कम उतार-चढ़ाव
के साथ: \omterm{def:g1:landscapes-around-us:landform}{पहाड़ी} की चोटी पर वह \omterm{def:g6:water-underground:spring}{कुएँ}~A के नीचे से ऊँचा है, पर ज़मीन उससे कहीं
ज़्यादा ऊपर उठती है, इसलिए पानी सतह से और नीचे होता है (काट पर, 110~मीटर से
भी ज़्यादा नीचे)।
\end{solution}

\begin{solution}{exo:g6:water-underground:10}
चिकनी मिट्टी के \omterm{def:g6:water-underground:porosity}{रंध्र} नन्हे हैं: पानी परतों से चिपका रहता है और चिकनी मिट्टी
में से \omterm{def:g6:water-underground:spring}{कुएँ} की ओर नहीं बह पाता। चिकनी मिट्टी \omterm{def:g6:water-underground:porosity}{सरंध्र} है पर \omterm{def:g6:water-underground:permeability}{अपारगम्य}, इसलिए उसमें
से केवल थोड़ा पानी (500 में से 20~लीटर) बाहर निकलता है।
\end{solution}

\begin{solution}{exo:g6:water-underground:11}
(क) $30.1 \div 0.26 \approx 116$: सौ गुना से भी ज़्यादा।
(ख) $100 - 69 = 31$: लगभग 31~\% बचता है, और उसमें से अकेला \omterm{def:g6:water-underground:groundwater}{भूजल} लगभग 30 है।
(ग) यह तरल \omterm{def:g2:water-on-earth:salt-fresh}{मीठे पानी} का अब तक का सबसे बड़ा भंडार है; बर्फ़ जमी हुई है और
ज़्यादातर बहुत दूर, ध्रुवीय क्षेत्रों में और पहाड़ों पर है। और \omterm{def:g6:water-underground:groundwater}{भूजल} हमारे
पैरों के नीचे लगभग हर जगह मिलता है।
\end{solution}

\begin{solution}{exo:g6:water-underground:12}
दोनों \omterm{def:g6:water-underground:spring}{कुएँ} उसी खड़िया \omterm{def:g6:water-underground:aquifer}{जलभृत} से पानी लेते हैं। शक्तिशाली पंप ने पानी उससे ज़्यादा
तेज़ी से निकाला जितनी तेज़ी से \omterm{def:g4:clouds-rain-snow:precipitation}{बारिश} खड़िया को फिर से भरती है, इसलिए चारों ओर \omterm{def:g6:water-underground:water-table}{भौम
जलस्तर} नीचे चला गया, दूसरे \omterm{def:g6:water-underground:spring}{कुएँ} की तली से भी नीचे, और वह \omterm{def:g6:water-underground:spring}{कुआँ} सूख गया। किसान
कम पंप करने पर सहमत हो सकते हैं, ख़ासकर गर्मियों में, एक ही \omterm{def:g6:water-underground:spring}{कुआँ} साझा कर सकते
हैं, या दूसरे \omterm{def:g6:water-underground:spring}{कुएँ} को और गहरा कर सकते हैं; सबसे बढ़कर, उन्हें उससे ज़्यादा पानी
नहीं लेना चाहिए जितना \omterm{def:g4:clouds-rain-snow:precipitation}{बारिश} लौटाती है।
\end{solution}

\begin{solution}{pb:g6:water-underground:1}
\textbf{1.} फ़ार्म खड़िया पर है; उसके नीचे शेल है। खड़िया \omterm{def:g6:water-underground:permeability}{पारगम्य} है, शेल
नहीं।

\textbf{2.} \omterm{def:g4:clouds-rain-snow:precipitation}{बारिश} खड़िया में रिसती है और शेल पर रुक जाती है। जहाँ खड़िया का आधार
ढलान पर बाहर निकलता है, वहाँ \omterm{def:g6:water-underground:water-table}{भौम जलस्तर} ज़मीन से मिलता है और पानी बाहर आता है:
यही \omterm{def:g6:water-underground:spring}{सोता} है।

\textbf{3.} $401 - 336 = 65$~मीटर ज़मीन के नीचे।

\textbf{4.} \omterm{def:g6:water-underground:spring}{कुएँ} की तली $401 - 100 = 301$~मीटर पर है, जो 314~मीटर पर स्थित
\omterm{def:g6:water-underground:water-table}{भौम जलस्तर} से नीचे है: \omterm{def:g6:water-underground:spring}{कुआँ} अब भी पानी तक पहुँचता है, $314 - 301 =
13$~मीटर से।

\textbf{5.} $100 \times 50 \times 2 = 10\,000$~घन मीटर।

\textbf{6.} $10\,000 \times 20 \div 100 = 2000$~घन मीटर।

\textbf{7.} $10\,000 \times 18 \div 100 = 1800$~घन मीटर बाहर निकलते हैं;
$2000 - 1800 = 200$~घन मीटर अटके रहते हैं।

\textbf{8.} $1800 \times 1000 = 1\,800\,000$~लीटर।

\textbf{9.} $18 \times 1000 = 18\,000$~लीटर रोज़।

\textbf{10.} $18 \times 7 = 126$~घन मीटर एक हफ़्ते में।

\textbf{11.} बरसात के मौसम की \omterm{def:g4:clouds-rain-snow:precipitation}{बारिश} खड़िया में रिसती है और हर साल उसे फिर से भर
देती है (जब तक फ़ार्म उससे ज़्यादा पंप न करें जितना \omterm{def:g4:clouds-rain-snow:precipitation}{बारिश} लौटाती है)।

\textbf{12.} $1800 \div 18 = $~\textbf{100 दिन}: बिना किसी \omterm{def:g4:clouds-rain-snow:precipitation}{बारिश} के तीन
महीने से थोड़ा ज़्यादा।
\end{solution}
